\section{Reproduction and archive map}

\subsection{Run the executable evidence}

No third-party Python packages are needed. From the archive's
\code{prototype} directory:
\par\noindent\begin{minipage}{\linewidth}
\begin{lstlisting}
python -S -m unittest discover -s tests -v
python -S replay.py
python -S run_experiments.py
python -S emit_lean.py
\end{lstlisting}
\end{minipage}\par

By default, the experiment generator writes to \code{reproduced-results} rather
than overwriting the retained \code{results}. The replay command defaults to
the retained results. The original generation command explicitly selected the
recorded output directory. Source, receipt, index, and order-ablation data are
deterministic for the recorded seed; elapsed time and environment metadata are
not expected to be identical across machines.

A small API example constructs and checks a Taylor certificate:
\par\noindent\begin{minipage}{\linewidth}
\begin{lstlisting}[language=Python]
from forge_analytic.core import taylor_remainder, encode
from forge_analytic.search import ladder
from forge_analytic.check import verify_goal

p = taylor_remainder(8)
request = {
    "polynomial": encode(p),
    "goal": {"kind": "positive_open_ray", "anchor": "0"},
}
receipt = ladder(p)
assert receipt is not None
assert verify_goal(request, receipt)
\end{lstlisting}
\end{minipage}\par

The same interface can be used with \code{make}, \code{ghost_search},
\code{tail}, and \code{compact_cover}; their docstrings state the exact input
and result conventions. These are algorithmic prototypes, not a call to a Lean
kernel through Python.

\subsection{Build the article and inspect the candidates}

The article uses ordinary pdfLaTeX and stock TeX packages:
\par\noindent\begin{minipage}{\linewidth}
\begin{lstlisting}
cd article
pdflatex -interaction=nonstopmode -halt-on-error forge-analytic-certificates.tex
pdflatex -interaction=nonstopmode -halt-on-error forge-analytic-certificates.tex
pdflatex -interaction=nonstopmode -halt-on-error forge-analytic-certificates.tex
\end{lstlisting}
\end{minipage}\par

The Lean directory includes its own status and intended compilation commands.
It does not include a vendored Mathlib checkout or a claim that an unchecked
candidate is a finished proof. The baseline reported by the audited Forge
snapshot is Lean \code{v4.34.0} and Mathlib
\pathline{1cf325a0cf67aca2b04d76b5380ff6a9e410aefa}~\cite{forge-readme}.

\subsection{Files worth opening first}

\begin{center}
\begin{tabularx}{\linewidth}{@{}p{0.38\linewidth}X@{}}
\toprule
Path & Purpose\\
\midrule
\path{article/} & Editable article, section sources, references, and rendered PDF\\
\path{prototype/forge_analytic/} & Exact algebra, search, enclosures, and separate replay checks\\
\path{prototype/tests/} & Unit tests, independent jet oracle, and corruption controls\\
\path{results/experiment-summary.json} & Counts, ablations, unknown controls, and environment\\
\path{results/problems.json} & Independently supplied canonical source requests\\
\path{results/certificates.json} & Recorded receipts linked to those requests\\
\path{lean/} & Candidate analytic bridge, four emitted replays, and explicit status\\
\path{docs/} & Repository audit, scope, and capability ledger\\
\bottomrule
\end{tabularx}
\end{center}

No original Forge proposal, third-party implementation, binary toolchain, or
font file is redistributed. The archive is an independent extension package;
it does not modify the user's repository.
